No equivariant assignment of classical partner values to the dying directions exists, in the minimal ascent or in any enlargement. Mechanism: the stabiliser of the null direction $[v]$ contains the parabolic null rotations, the candidate partners form an affine space on which those rotations act by translation, and a translation of an affine space has no fixed point; cohomologically, the obstruction is the tautological class of the unipotent stabiliser, nonzero in every mixed-signature dimension above two. \emph{Provenance:} the element-level no-go is \cite[Thm.~4.2]{ButlerT1}; the totality of the transported division is \cite[Thm.~4.10]{ButlerT1}; the cohomological form, with the torsor made exact modulo the left ideal and the class identified as $\theta \in H^1(N_v; U) = \operatorname{Hom}(N_v, U)$, is \cite[Thm.~4.8]{ButlerT1}. The dimension-two exception, where the class vanishes, is exactly Lemma \claimref{singularity.pseudo-inverse-identities-hyperbolic-atom}'s self-completeness. \emph{Scope:} these theorems are proved where the ambient form is nondegenerate and the dying direction has a hyperbolic partner, which is the situation at every point of the invertible region, whose limits toward $E$ the no-go therefore governs; at $E$ itself the partner is exactly what is missing (Proposition \claimref{singularity.type-n-only-type-n}), so the import does not apply there verbatim. The boundary case is supplied by the no-value theorem at the wall, \cite[Thm.~10.4]{ButlerT4}, reproduced verbatim as Fact \claimref{singularity.imported-no-equivariant-value-wall}: no equivariant classical value exists at $E$ either, the completion's own transversal-independence carrying the equivariance. The claim that no classical value can repair the dying directions therefore holds on the approach by \cite[Thm.~4.2]{ButlerT1} and at the boundary by Fact \claimref{singularity.imported-no-equivariant-value-wall}.
